\documentclass[10pt]{article}
\usepackage[margin=0.65in]{geometry}
\usepackage{amsmath}
\usepackage{booktabs}
\usepackage{array}
\usepackage[T1]{fontenc}

\setlength{\parindent}{0pt}
\setlength{\parskip}{0.35em}

\begin{document}

\begin{center}
{\Large Two-Track DCA$_{xy}$ Comparison}\\[0.25em]
{\small FastSim closest-point residual vs. cross-product projection}
\end{center}

\section*{Definitions}

For track PCA positions $\mathbf{r}_i$, $\mathbf{r}_j$ and momentum directions
$\mathbf{u} = \mathbf{p}_i / |\mathbf{p}_i|$,
$\mathbf{v} = \mathbf{p}_j / |\mathbf{p}_j|$, define
\[
  \mathbf{n} = \frac{\mathbf{u} \times \mathbf{v}}
                    {|\mathbf{u} \times \mathbf{v}|},
  \qquad
  \Delta\mathbf{r} = \mathbf{r}_j - \mathbf{r}_i .
\]

The cross-product method discussed here is
\[
  D^{\mathrm{cross}}_{xy}
    = \mathbf{n}\cdot(\Delta x,\Delta y,0).
\]
This is a projection of the transverse PCA separation onto the normal vector
of the two track directions.

The current FastSim method first solves the two closest points
$\mathbf{c}_i$ and $\mathbf{c}_j$ on the full 3D track lines, then computes
\[
  \boldsymbol{\rho} = \mathbf{c}_i - \mathbf{c}_j,
  \qquad
  D^{\mathrm{FS}}_{xy}
    = s\,\sqrt{\rho_x^2+\rho_y^2},
\]
where the sign $s$ is set from
\[
  s = \mathrm{sign}\left[
    \left((\mathbf{p}_{T,i}+\mathbf{p}_{T,j})
    \times(\mathbf{r}_{i}-\mathbf{r}_{j})\right)_z
  \right].
\]

\section*{Main Result}

These are not the same observable. The FastSim value is the transverse
magnitude of the closest-point residual. The cross-product value is a
normal-vector projection after forcing the PCA separation's $z$ component
to zero. They can have similar overall scale, but they are not related by a
single event-by-event scale factor.

Signed values should not be compared directly without first harmonizing the
sign convention. The table below compares absolute magnitudes.

\section*{Numerical Check}

The following check used the default FastSim configuration, seed 20260509,
20,000 generated events, and 12,919 selected Si-Si pairs. All distances are
in cm.

\begin{center}
\begin{tabular}{lrrrrrr}
\toprule
Quantity & Median & 68\% & 90\% & 95\% & 99\% & Max \\
\midrule
$|D^{\mathrm{FS}}_{xy}|$       & 0.00500 & 0.00819 & 0.01683 & 0.02260 & 0.03636 & 0.08340 \\
$|D^{\mathrm{cross}}_{xy}|$    & 0.00419 & 0.00710 & 0.01578 & 0.02174 & 0.03611 & 0.08151 \\
$\left||D^{\mathrm{cross}}_{xy}|-|D^{\mathrm{FS}}_{xy}|\right|$
                                  & 0.00240 & 0.00400 & 0.00844 & 0.01159 & 0.01925 & 0.03683 \\
\bottomrule
\end{tabular}
\end{center}

The ratio $|D^{\mathrm{cross}}_{xy}|/|D^{\mathrm{FS}}_{xy}|$ had median 0.91.
Its 5--95\% interval was 0.087--6.64, with a 99\% value of 32.7. The large
upper tail appears when the FastSim denominator is close to zero.

\section*{Interpretation}

For resolution-width comparisons, the two definitions may look broadly similar
because both are controlled by the same PCA and angular scales. For event-by-event
comparisons, DCA cuts, or tail studies, the difference can be substantial. A few
$10^{-3}$ cm is typical in the default toy sample, while the tail reaches a few
$10^{-2}$ cm.

\end{document}
